\documentclass[11pt,a4paper]{article}
\usepackage[T1]{fontenc}
\usepackage{lmodern}
\usepackage{amsmath,amssymb,amsthm,mathtools,booktabs,longtable,graphicx,microtype}
\usepackage[margin=24mm,headheight=15pt]{geometry}
\usepackage{hyperref}
\usepackage{fancyhdr}
\hypersetup{colorlinks=true,urlcolor=blue!55!black,linkcolor=blue!55!black,citecolor=blue!55!black,pdftitle={Quantum Action, Topology and Identifiability}}
\usepackage{xcolor}
\newtheorem{proposition}{Proposition}[section]
\newtheorem{corollary}[proposition]{Corollary}
\newtheorem{remark}[proposition]{Remark}
\newcommand{\h}{\hbar}
\newcommand{\R}{\mathbb R}
\newcommand{\C}{\mathbb C}
\newcommand{\Z}{\mathbb Z}
\newcommand{\T}{\mathbb T}
\newcommand{\tr}{\operatorname{tr}}
\newcommand{\rank}{\operatorname{rank}}
\newcommand{\KL}{D_{\mathrm{KL}}}
\newcommand{\Span}{\operatorname{span}}
\pagestyle{fancy}
\fancyhf{}\fancyhead[L]{Research memorandum --- mathematical and physical status separated}\fancyhead[R]{17 October 2026}\fancyfoot[C]{\thepage}
\begin{document}
\begin{center}
{\Large\bf Quantum Action, Topology, and Identifiability:}\\[2mm]
{\large\bf A Critical Search for a Deeper Origin of Planck's Constant}\\[5mm]
{\normalsize Independent mathematical synthesis and computational research memorandum}\\[2mm]
{\small Working manuscript; scientific authorship and publication approval pending}\\[1mm]
{\small 17 October 2026 --- Not peer reviewed}
\end{center}
\begin{abstract}
We ask what the structure of quantum theory, symplectic geometry, topology, phase information, and gravitational models actually explains about Planck's constant. We separate three inequivalent tasks: deriving quantum noncommutativity, selecting a physical quantum-action scale, and explaining the decimal value fixed by SI units. We derive precise obstructions to the stronger task under common assumptions. The regular Weyl canonical commutation relations possess models for all positive $\h$; canonical cotangent-bundle symplectic forms on $T^*(S^3/\Gamma)$ are exact, so their prequantization does not discretize $\h$; discrete flat line-bundle holonomies and harmonic selection rules label topological sectors without calibrating action. We further establish an exact amplitude--$\h$ nonidentifiability for covariance models of the form $N+\h A K_\Gamma$. Conversely, the Weil condition on a \emph{fixed non-exact symplectic form} and independently normalized interference action show what sorts of additional hypotheses can constrain a quantum scale. A published proposal to derive $\h$ from spinorial phase winding is evaluated by dimensional and representation-theoretic counterchecks. We provide a reproducible code suite with 9,886 finite-dimensional/numerical checks. The results are structural and operational; they do not constitute a derivation of the observed numerical value of $\h$ or a new fundamental physical law.
\end{abstract}
\tableofcontents
\section{Scientific question, units, and evidence status}
The reduced Planck constant $\h=h/(2\pi)$ has physical dimensions of action. In the revised SI the \emph{numerical value} of $h$ is fixed exactly:
\begin{equation}\label{eq:si}
 h=6.62607015\times10^{-34}\ \mathrm{J\,s},\qquad
 \h\simeq 1.054571817646\times10^{-34}\ \mathrm{J\,s}.
\end{equation}
The SI value is an operational definition of the kilogram, not a microscopic explanation of quantization \cite{NIST-SI,NIST-kibble}. Its digits do not form a unit-independent prediction. The scientifically discriminating questions concern the existence of a common action scale, how it compares to independently defined physical quantities, and whether a dynamical theory predicts \emph{dimensionless} ratios \cite{DuffOkunVeneziano,CODATA}.

We distinguish \textbf{established mathematical claims} under explicitly stated assumptions from \textbf{proposed dynamical hypotheses}. Reproducing a conventional identity in new notation is not independent proof of a physical law. The user's META--CAUSAL/$S^3$ monograph, v244, includes conformal-vacuum covariances proportional to $\h$ and topological mode projectors, but explicitly chooses canonical commutators $[q,p]=i\h$ and in other sections uses natural units $\h=c=1$. It therefore cannot presently be read as deriving $\h$.

\section{Continuous families of quantum kinematics}
\begin{proposition}[Heisenberg--Weyl representations at arbitrary quantum scale]\label{pr:weyl}
For every real $\h>0$ there is a regular irreducible representation on $L^2(\R)$ of
\begin{equation}\label{eq:CCR}
  (X\psi)(x)=x\psi(x),\qquad (P_{\h}\psi)(x)=-i\h\frac{d\psi}{dx}(x),\qquad
  [X,P_{\h}]=i\h I,
\end{equation}
on their common Schwartz domain. In Weyl form, with $U(a)=e^{-iaP_\h/\h}$ and $V(b)=e^{ibX/\h}$,
\begin{equation}\label{eq:weyl}
 U(a)V(b)=e^{-iab/\h}V(b)U(a).
\end{equation}
\end{proposition}
\begin{proof} $U(a)\psi(x)=\psi(x-a)$ and $V(b)\psi(x)=e^{ibx/\h}\psi(x)$. Substitution gives the displayed Weyl relation. The representation is regular; irreducibility follows by the standard Schr\"odinger/Weyl argument. The Stone--von Neumann theorem gives uniqueness only \emph{at a specified central character}, not a selection of that character \cite{Hall}.\end{proof}
\begin{corollary}[Kinematics alone does not select a numerical action scale]
The canonical algebra, continuity and irreducibility cannot separately determine a particular positive $\h$. Distinct choices are not necessarily physically equivalent when classical masses, charges and calibrated momenta are held fixed; however those scales are extra empirical input, not a consequence of the Weyl algebra.
\end{corollary}
\begin{proposition}[Finite-dimensional trace obstruction]
There exist no finite-dimensional matrices $X,P$ with $[X,P]=i\h I$ for $\h\ne0$.
\end{proposition}
\begin{proof} $\tr[X,P]=0$, whereas $\tr(i\h I)=i\h d\ne0$. Finite-dimensional clock--shift operators instead satisfy exponentiated root-of-unity relations, with a different algebra.\end{proof}
\begin{remark}[The uncertainty principle does not numerically fix $\h$]
For any $\h>0$ the Gaussian oscillator state with
\begin{equation}
 V_{\h}=\begin{pmatrix}\frac{\h}{2m\omega}&0\\0&\frac{\h m\omega}{2}\end{pmatrix}
 \quad\text{satisfies}\quad V_{\h}+\frac{i\h}2\begin{pmatrix}0&1\\-1&0\end{pmatrix}\succeq0,
\qquad \det V_{\h}=\frac{\h^2}4.
\end{equation}
The inequality constrains a \emph{ratio} of measured uncertainty and the independently calibrated action parameter; it does not pick the latter by itself.
\end{remark}

\section{Symplectic geometry: when topology constrains quantization}
Let $Q$ be a configuration manifold and $(T^*Q,\omega)$ its canonical phase space. Denote by $\lambda$ the tautological Liouville one-form and take $\omega=d\lambda$ up to the conventional sign.
\begin{proposition}[No Weil-period restriction for the canonical cotangent bundle]\label{pr:cotangent}
For every $\h>0$, the canonical symplectic form on $T^*Q$ satisfies the prequantum integrality condition because its de Rham cohomology class vanishes:
\begin{equation}\label{eq:exact}
 \frac{[\omega]}{2\pi\h}=0\ \text{in }H^2_{\mathrm{dR}}(T^*Q), \qquad
 \frac{1}{2\pi\h}\int_\Sigma\omega=0
\end{equation}
for each closed two-cycle $\Sigma$. In particular this holds for $Q=S^3$ and $Q=L(p,q)$.
\end{proposition}
\begin{proof} The symplectic form $\omega$ is exact. Stokes's theorem makes its integral over any closed oriented two-cycle vanish. A globally trivial prequantum line with connection proportional to $\lambda/\h$ realizes its curvature for each positive $\h$. This is a statement about existence of \emph{prequantization}, not canonical choice of a quantum state or polarization \cite{GeoQuant}.\end{proof}
\begin{proposition}[Fixed non-exact symplectic periods can restrict ratios]\label{pr:weil}
For a \emph{specified} symplectic manifold $(M,\Omega)$, a prequantum line bundle requires $[\Omega/(2\pi\h)]$ to be an integral cohomology class. If $M=S^2$ with total symplectic area $A=\int_{S^2}\Omega>0$, then the condition is
\begin{equation}\label{eq:weil}
 \frac{A}{2\pi\h}\in\Z_{>0}\,,\qquad \h=\frac{A}{2\pi k},\ k\in\Z_{>0}.
\end{equation}
If $A$ is independently specified in physical units and $k$ is fixed by additional physical information, the relation constrains $\h$; topology alone supplies only $k$.
\end{proposition}
\begin{proof} Integrality of the first Chern class is the Weil condition, evaluated on the fundamental 2-cycle of the sphere \cite{GeoQuant}. Both $A$ and $k$ are necessary to obtain a value.\end{proof}
\begin{remark}[Why the three-sphere itself is not a symplectic phase space]
A nondegenerate alternating two-form exists only in even dimensions. Thus the odd-dimensional spatial manifold $S^3$ cannot carry a symplectic form. The relevant Hamiltonian phase space, e.g. $T^*S^3$, has even dimension six and canonical exact $\omega$. A geometric $S^3$ radius does not by itself supply an action period.
\end{remark}

\section{Lens spaces: discrete topological sectors without an action quantum}
The lens space $L(p,q)$ is $S^3$ modulo the free generator $(z_1,z_2)\mapsto(e^{2\pi i/p}z_1,e^{2\pi iq/p}z_2)$ with $\gcd(p,q)=1$.
\begin{proposition}[Discrete flat sectors independent of $\h$]\label{pr:flat}
For $p>1$,
\begin{equation}
 \pi_1(L(p,q))=\Z_p,\quad H^2(L(p,q);\Z)\cong\Z_p,\quad
 H^2_{\mathrm{dR}}(L(p,q))=0,
\end{equation}
and flat unitary characters are
\begin{equation}\label{eq:holonomy}
 \chi_a(1)=e^{2\pi ia/p},\quad a=0,\ldots,p-1.
\end{equation}
Thus there are $p$ scalar character sectors, whose holonomies are dimensionless roots of unity. These sectors coexist with every positive $\h$; they do not select its magnitude.
\end{proposition}
\begin{proof} The standard CW complex for $L(p,q)$ has one cell in dimensions $0,1,2,3$ and cellular boundary $\partial_2=p$, giving $H_1=\Z_p$ and $H_2=0$. The universal coefficient theorem gives $H^2\cong\operatorname{Ext}(H_1,\Z)=\Z_p$. Torsion vanishes after passing to real cohomology. A flat $U(1)$ line is classified by its holonomy character $\operatorname{Hom}(\pi_1,U(1))$, giving the $p$ roots in (\ref{eq:holonomy}) \cite{ImboSudarshan}. A prequantum connection of canonical curvature may be tensored with these flat lines without changing its curvature or $\h$.\end{proof}
\begin{proposition}[Character projectors and harmonic multiplicities]\label{pr:projector}
Let $\rho_n(g)$ be the action of the lens generator on the degree-$n$ spherical harmonic space $\mathcal H_n$. The orthogonal projector on character sector $a$ is
\begin{equation}\label{eq:proj}
 P_{n,a}=\frac1p\sum_{k=0}^{p-1}e^{-2\pi i ak/p}\rho_n(g)^k,
 \quad P_{n,a}^*=P_{n,a}=P_{n,a}^2.
\end{equation}
For scalar harmonic polynomials obtained from $\C^2\simeq\R^4$, let $M_{n,a}$ count monomials $z_1^{u}z_2^{v}\bar z_1^{w}\bar z_2^{x}$ of degree $n$ satisfying $(u-w)+q(v-x)\equiv a\pmod p$. Then
\begin{equation}\label{eq:count}
 \dim\mathcal H_{n,a}=M_{n,a}-M_{n-2,a},\qquad M_{m,a}=0\ (m<0).
\end{equation}
These formulae depend on $(p,q,n,a)$ but not on $\h$. In particular for $n=2$, the invariant multiplicities are $3$ on $L(5,1)$ and $1$ on $L(5,2)$.
\end{proposition}
\begin{proof} Character orthogonality in $\Z_p$ proves idempotence. Every homogeneous polynomial is a direct sum of a harmonic polynomial and $r^2$ times a polynomial of degree $n-2$; $r^2=|z_1|^2+|z_2|^2$ is group-invariant. Taking character spaces gives (\ref{eq:count}). Counting the monomials of degree 2 proves the example.\end{proof}
\begin{figure}[ht]\centering\includegraphics[width=.85\linewidth]{figures/lens_scalar_multiplicities.pdf}
\caption{Exact invariant scalar harmonic multiplicities in two lens spaces with the same fundamental-group order. A topology-sensitive spectral distinction does not constitute a prediction of the quantum action constant.}\end{figure}

\section{The cosmological covariance identifiability barrier}
The META--CAUSAL/S$^3$ research program includes quantum vacuum correlators, group projectors, finite observation operators and nuisance-profiled information measures. All are compatible with the following exact statistical obstruction.
\begin{proposition}[Amplitude--quantum-action degeneracy]\label{pr:degeneracy}
Suppose data are modeled as a zero-mean Gaussian with positive definite covariance
\begin{equation}\label{eq:cov}
 C_\Gamma(\h,A,\theta)=N(\theta)+\h A K_\Gamma(\theta),\quad \h,A>0,
\end{equation}
where $K_\Gamma\succeq0$ is fixed for each topology and $A$ is an unconstrained positive source amplitude. For all $s>0$, the transformation
\begin{equation}\label{eq:gauge}
 (\h,A)\mapsto(s\h,A/s)
\end{equation}
leaves the entire Gaussian probability law invariant. Consequently $\h$ is not globally identifiable from these observations, even with arbitrarily many samples from the same covariance family. At a regular fiducial point with nonzero signal, the Fisher information for $\h$ after profiling $A$ equals zero.
\end{proposition}
\begin{proof} Substitution into (\ref{eq:cov}) proves exact covariance equality. The two score directions satisfy $\partial_{\h}C=A K_\Gamma$ and $\partial_A C=\h K_\Gamma$, so they are proportional. The covariance Fisher metric $F_{ij}=\frac12\tr(C^{-1}\partial_i C\,C^{-1}\partial_j C)$ has a rank-one block in $(\h,A)$ and its Schur complement with respect to $A$ is zero. Profiling additional nuisance parameters cannot recover a direction already annihilated by exact model equivalence.\end{proof}
\begin{corollary}[Topology itself cannot self-calibrate $\h$]
Knowing $P_n^\Gamma$, all scalar and tensor harmonic multiplicities, and even the full transfer operator does not remove the degeneracy of Proposition~\ref{pr:degeneracy} while the signal contains only the product $\h A$ and the other inputs are not independently calibrated. Observing more $n$-shells does not break that \emph{common} scale symmetry.
\end{corollary}
\begin{figure}[ht]\centering\includegraphics[width=.85\linewidth]{figures/hbar_amplitude_degeneracy.pdf}
\caption{Different choices of quantum-action scale $\h$ and unknown source amplitude $A$ yielding identical covariances on each hyperbola. Units are illustrative.}\end{figure}
\begin{proposition}[Minimal independent calibration criterion]\label{pr:calibration}
Let $u=\log\h$ and $v=\log A$. Suppose two independently measured positive observables satisfy locally
\begin{equation}
 \log O_1=u+v+c_1,\qquad \log O_2=r u+s v+c_2,
\end{equation}
with known $c_1,c_2,r,s$ and nonsingular measurement noise. They locally identify $(u,v)$ if and only if $s\ne r$, because
\begin{equation}\label{eq:rank}
 \det\begin{pmatrix}1&1\\r&s\end{pmatrix}=s-r.
\end{equation}
This is a necessary-and-sufficient local rank criterion for this two-parameter calibrated model, not a derivation of the input physical constants.
\end{proposition}
\begin{proof} The Jacobian of the logarithmic observation map is constant, and the inverse-function theorem applies exactly when its determinant is nonzero.\end{proof}
\begin{remark}
An externally calibrated interferometric action difference $\Delta S$ gives $\Delta\varphi=\Delta S/\h$ modulo $2\pi$. Its normalization is different from an unknown variance amplitude and can, in principle, break a covariance scale degeneracy. A phase-wrap ambiguity and empirical uncertainties must still be controlled. Such a measurement calibrates $\h$; it does not prove an ontological origin.
\end{remark}
\begin{figure}[ht]\centering\includegraphics[width=.75\linewidth]{figures/finite_noise_fisher.pdf}
\caption{For a fixed identifiable direction, increased noise suppresses Fisher information smoothly, not at a universal finite threshold. Note that the exact zero Fisher of Proposition~\ref{pr:degeneracy} is a different structural phenomenon.}\end{figure}

\section{Topological phase winding: an audit of proposed derivations}
A July 2025 unreviewed preprint claimed to derive $\h$ from phase winding and a minimal spinorial action threshold \cite{Altunin}. The topic is relevant, but its dimensional step must be scrutinized, rather than accepted as a physical theorem.
\begin{proposition}[Winding does not set a dimensionful action]\label{pr:winding}
For an ordinary single-valued complex scalar $\Psi=\rho e^{i\theta}$ with nonzero amplitude on a closed loop, one has
\begin{equation}\label{eq:winding}
 \frac1{2\pi}\oint d\theta=n\in\Z.
\end{equation}
The integer $n$ alone cannot fix any positive number with dimensions of action. In the elementary phase-gradient model on a circle of circumference $L$, with stiffness $\kappa>0$, the minimum gradient energy in sector $n$ is
\begin{equation}\label{eq:stiffness}
 E_n=\kappa\int_0^L (\partial_x\theta)^2dx
 =\frac{4\pi^2\kappa n^2}{L}.
\end{equation}
The energy scale varies continuously with the \emph{independent} stiffness $\kappa$ and size $L$ while the integer winding remains unchanged. Converting energy to action requires further independently specified time data.
\end{proposition}
\begin{proof} A continuous map $S^1\to S^1$ has integer degree. For fixed phase increment $2\pi n$, Cauchy--Schwarz yields $\int_0^L(\theta')^2dx\ge(2\pi n)^2/L$, with equality for constant derivative. Rescaling $\kappa$ changes the minimum continuously, leaving winding unchanged.\end{proof}
\begin{remark}[Spinor monodromy is an extra structure]
A sign change $\Psi(2\pi)=-\Psi(0)$ is not a closed loop in the nonzero \emph{ordinary complex scalar} field itself. It can be a consistent boundary condition for an appropriate spinor or section with transition functions, but such a spin structure and representation must be specified independently. A nontrivial sign after $2\pi$ rotation does not by itself fix any energy or action unit.
\end{remark}
\begin{remark}[Quantum coherence has no general hard cutoff $|\Delta S|\ge\h$]
Two-path interference includes the factor $\cos(\Delta S/\h)$. For $0<|\Delta S|\ll\h$, the phase is small but the interference is not generically destroyed. A claim that all physically distinguishable structures vanish when an action increment is below $\h$ requires explicit dynamics and contradicts the generic coherence of arbitrarily small phase differences under ideal repeated measurements. Small effects can be difficult to detect, but are not logically absent.
\end{remark}

\section{Quantum gravity, $S^3$, and possible ultraviolet scales}
In common loop-quantum-gravity formulations, area eigenvalues are of the schematic form
\begin{equation}\label{eq:lqg}
 A_{\{j_e\}}=8\pi\gamma\,\frac{G\h}{c^3}\sum_{e\cap\Sigma}\sqrt{j_e(j_e+1)}.
\end{equation}
Here $\gamma$ is the dimensionless Barbero--Immirzi parameter; the equation already assumes $\h$ and $G$ and cannot determine $\h$ solely from spin labels. The quantization ambiguity has been documented in the foundational literature \cite{RovelliThiemann}. Likewise, the Planck length $\ell_P^2=G\h/c^3$ \emph{defines} a combination of inputs, rather than predicting one from the others.

A physically motivated emergent-action program would have to produce some dimensionless ratio such as
\begin{equation}\label{eq:alphag}
 \alpha_G(m)=\frac{Gm^2}{\h c},\qquad
 \frac{\h}{mcR},\qquad
 \frac{\Delta S_{\mathrm{cal}}}{\h}
\end{equation}
from a single dynamical theory and independently defined observables. Invoking a hypothetical preferred cosmic radius $R$ or a topological integer $p$ adds no physical mass or action scale. No selection is proven merely by writing $\h=F(G,c,R,p,\ldots)$ when the numerical normalization or a new dimensional parameter is freely chosen.

\subsection{The most stringent constructive research target}
A genuine first-principles result would require all of the following:
\begin{enumerate}
\item A state space and dynamics defined independently of the quantity to be derived, including dimensions and normalization.
\item A theorem yielding a stable, universal quantum-action scale in many unrelated systems, not only a single mode or interference device.
\item A new, calibrated, dimensionless numerical prediction that is not already an assumption in the model.
\item A deviation or correlation measurable independently of calibration conventions and established quantum electrodynamics.
\item Consistency of the model with the observed common quantum scale and existing constraints on its variation.
\end{enumerate}
The preceding propositions localize a precise obstruction that any such model must overcome: continuous quantization freedom, exact phase-space symplectic forms, topological torsion sectors, and amplitude nonidentifiability.

\section{Reproducibility protocol and boundaries}
The executable \texttt{verify\_hbar\_foundations.py} tests 9,886 finite numerical/algebraic cases: Weyl commutation at varying $\h$; the uncertainty matrix; finite-dimensional CCR trace anomaly; lens-space harmonic counting versus finite character averaging; hermitian idempotent projectors; scale-invariant covariance and zero profile Fisher; quadratic KL expansion; winding energy; and flux/holonomy quantization. The separate plotting script produces all vector figures. The tests \emph{do not} check a physical unified action theory, LQG dynamics or observational data. The standalone source files from META--CAUSAL v244 are preserved separately; no unreviewed new theorem is inserted into the original PRD monograph.

\begin{center}
\fbox{\begin{minipage}{.91\linewidth}
\textbf{Final status:} The structure and operational role of $\h$ admit rigorous mathematical explanation. A unique numerical derivation from $S^3$, finite lens-space topology, phase winding, Fibonacci structure, symplectic integrality or generic information theory has \emph{not} been established. The positive, testable achievement here is to prove exact underdetermination mechanisms and state the missing independent calibration condition. The manuscript claims neither priority over the cited mathematics nor a newly discovered fundamental physical law.
\end{minipage}}
\end{center}

\begin{thebibliography}{99}
\bibitem{NIST-SI} National Institute of Standards and Technology, \emph{Definitions of SI Base Units}, updated 2025, \url{https://www.nist.gov/si-redefinition/definitions-si-base-units}.
\bibitem{NIST-kibble} S. Schlamminger and D. El Haddad, \emph{The Kibble balance and the kilogram}, 2019, \url{https://www.nist.gov/publications/kibble-balance-and-kilogram}.
\bibitem{CODATA} P. J. Mohr, D. B. Newell, B. N. Taylor and E. Tiesinga, \emph{CODATA recommended values of the fundamental physical constants: 2022}, Rev. Mod. Phys. 97, 025002 (2025), \url{https://doi.org/10.1103/RevModPhys.97.025002}.
\bibitem{DuffOkunVeneziano} M. J. Duff, L. B. Okun and G. Veneziano, \emph{Trialogue on the number of fundamental constants}, JHEP 03 (2002) 023, \url{https://doi.org/10.1088/1126-6708/2002/03/023}.
\bibitem{Hall} B. C. Hall, \emph{Quantum Theory for Mathematicians}, Springer (2013), Chs. 13 and 22; Weyl canonical commutation relations and geometric quantization.
\bibitem{GeoQuant} \emph{The regularity and products in contact geometry}, Ann. Mat. Pura Appl. (2025), discussion of the Weil integrality condition, \url{https://doi.org/10.1007/s10231-025-01631-7}.
\bibitem{ImboSudarshan} T. D. Imbo and E. C. G. Sudarshan, \emph{Inequivalent quantizations and fundamentally perfect spaces}, Phys. Rev. Lett. 60, 481 (1988), \url{https://doi.org/10.1103/PhysRevLett.60.481}.
\bibitem{RovelliThiemann} C. Rovelli and T. Thiemann, \emph{Immirzi parameter in quantum general relativity}, Phys. Rev. D 57, 1009 (1998), \url{https://doi.org/10.1103/PhysRevD.57.1009}.
\bibitem{Altunin} M. Altunin, \emph{Quantization Without Postulates: Deriving $\h$ from Phase Winding}, Preprints.org (2025), \emph{not peer reviewed}, \url{https://doi.org/10.20944/preprints202507.2668.v1}.
\bibitem{COMPACT} COMPACT Collaboration, \emph{CMB-for-S3: correlations for $S^3$ and lens spaces}, \url{https://github.com/CompactCollaboration/CMB-for-S3}.
\end{thebibliography}
\end{document}
